\documentclass[11pt, a4paper]{article}
\usepackage{amsmath}

\begin{document}
 
\section*{Calcolo del Rotore in Coordinate Curvilinee Ortogonali}

\subsection*{1. Definizione delle coordinate}
Identifichiamo le tre variabili indipendenti del sistema di coordinate ortogonali:
\begin{equation*}
    (q_1, q_2, q_3)
\end{equation*}

\subsection*{2. Calcolo dei fattori di scala}
Per ogni coordinata, definiamo il fattore di scala (o coefficiente di Lamé) $h_i$, che rappresenta il rapporto tra lo spostamento fisico infinitesimo $ds_i$ e la variazione della coordinata $dq_i$:
\begin{equation*}
    h_i = \frac{ds_i}{dq_i} \quad \text{per } i=1,2,3
\end{equation*}

\subsection*{3. Scomposizione del campo vettoriale}
Esprimiamo il campo vettoriale $\vec{F}$ lungo i versori ortogonali normalizzati $(\hat{u}_1, \hat{u}_2, \hat{u}_3)$:
\begin{equation*}
    \vec{F} = F_1(q_1, q_2, q_3) \hat{u}_1 + F_2(q_1, q_2, q_3) \hat{u}_2 + F_3(q_1, q_2, q_3) \hat{u}_3
\end{equation*}

\subsection*{4. Costruzione della matrice del rotore}
Impostiamo l'equazione del rotore utilizzando il fattore di normalizzazione esterno (il reciproco del determinante dello Jacobiano) e la matrice formale:
\begin{equation*}
    \vec{\nabla} \times \vec{F} = \frac{1}{h_1 h_2 h_3} 
    \begin{vmatrix} 
        h_1 \hat{u}_1 & h_2 \hat{u}_2 & h_3 \hat{u}_3 \\ 
        \frac{\partial}{\partial q_1} & \frac{\partial}{\partial q_2} & \frac{\partial}{\partial q_3} \\ 
        h_1 F_1 & h_2 F_2 & h_3 F_3 
    \end{vmatrix}
\end{equation*}

\subsection*{5. Sviluppo del determinante}
Risolvendo il determinante lungo la prima riga e semplificando i fattori esterni con quelli dei versori estratti, otteniamo l'espressione analitica finale per il rotore:
\begin{align*}
    \vec{\nabla} \times \vec{F} = 
    &\frac{1}{h_2 h_3} \left( \frac{\partial (h_3 F_3)}{\partial q_2} - \frac{\partial (h_2 F_2)}{\partial q_3} \right) \hat{u}_1 \\
    + &\frac{1}{h_1 h_3} \left( \frac{\partial (h_1 F_1)}{\partial q_3} - \frac{\partial (h_3 F_3)}{\partial q_1} \right) \hat{u}_2 \\
    + &\frac{1}{h_1 h_2} \left( \frac{\partial (h_2 F_2)}{\partial q_1} - \frac{\partial (h_1 F_1)}{\partial q_2} \right) \hat{u}_3
\end{align*}

\end{document}